\documentclass[11pt,a4paper]{article}
\usepackage[T1]{fontenc}
\usepackage[utf8]{inputenc}
\usepackage[polish]{babel}
\usepackage{amsmath,amssymb,amsthm,geometry,hyperref}
\geometry{margin=25mm}
\newtheorem{theorem}{Twierdzenie}
\newtheorem{corollary}{Wniosek}
\newcommand{\Tr}{\operatorname{Tr}}
\newcommand{\op}{\operatorname{op}}
\newcommand{\Rea}{\operatorname{Re}}
\newcommand{\id}{\operatorname{id}}
\title{FIN: korelacje mikroskopowe i koszt kwantowego odtworzenia jądra}
\author{}
\date{}
\begin{document}
\maketitle

\begin{abstract}
Naturalne dwuciałowe podniesienie rzutowanego źródła FIN ma bardzo
restrykcyjną strukturę: na dowolnym skończonym połączonym grafie,
nawet z dodatkowymi hamiltonianami lokalnymi, jego jedynym
pełnorządowym stanem iloczynowym w równowadze jest stan maksymalnie
mieszany. Niezerowy strict nie może więc być takim nieskorelowanym
stanem mikroskopowym. Jawne skorelowane stany o tym samym marginalnym
strict istnieją, lecz mogą tylko przechowywać jądro w paśmie płaskim,
bez odtwarzania jego propagacji. Osobno skonstruowano zgodny z QM
protokół odtwarzający kanał strict z kopii stanu programu. Ma on
ścisłą górną granicę błędu i nieusuwalny w tej architekturze składnik
błędu rzędu $1/N$. Źródło danych jądra pozostaje w programie.
Wyniki ustanawiają konkretne warunki dla dalszego modelu fizycznego;
nie są wyprowadzeniem FIN jako teorii fundamentalnej.
\end{abstract}

\section{Przedmiot twierdzeń i granice interpretacji}

Rozważany jest rzeczywisty rzut źródłowy
\begin{equation}\label{eq:source}
 T(\rho)=\Pi\Rea\rho,
\end{equation}
gdzie $\Pi$ usuwa przekątną. W proponowanym sprzężeniu stanu i jądra
$\dot K=\eta(T(\rho)-\gamma K)$, a granica szybkiego śledzenia daje
$\dot\rho=(i/\gamma)[T(\rho),\rho]$.
Jest to granica określonego równania, nie założenie, że każde
samouczenie musi mieć tę postać.
\footnote{Źródła repozytoryjne: \texttt{The FIN Kernel as an Unknown Mat.md},
§4.1; \texttt{fin\_projected\_learning/REPORT.md}, §28 (ST8648).
Wcześniejszy raport \texttt{FIN\_Adaptive\_Feedback\_Operational\_Obstruction.tex}
oddziela oryginalne równanie czyste od jego rozszerzenia mieszanego.}

Wprowadzamy jawnie dwa dodatkowe modele: skończone oddziaływanie
dwuciałowe oraz protokół kolejnych zderzeń z kopiami stanu programu.
Ich poprawność kwantowa nie oznacza, że prawo składania, dopływ kopii
lub zegar zostały już wyprowadzone z FIN. Liczba $n=12$ oznacza
wymiar pojedynczego nośnika, a $N$ liczbę składników lub kopii;
nie utożsamiamy $N$ z wymiarem fraktalnym, liczbą cząstek fizycznych
ani wykładnikiem hierarchii grawitacyjnej.

Punktem odniesienia pozostaje dokładnie zadany strict na cyklu:
\begin{equation}\label{eq:strict}
 W_{ij}=\frac{\cos(\omega d_{ij}+\phi)}{1+d_{ij}^{9/5}},\quad
 W_{ii}=0,\quad \omega=\frac{743}{4000},\quad\phi=\frac{650}{4000}.
\end{equation}
Jego sumę wiersza oznaczamy przez $s$, a $\ell=-\lambda_{\min}(W)>0$.
Wszystkie czasy i parametry są bezwymiarowe. Zasadnicze wyniki
o komutantach nie zależą od wartości tych parametrów.

\section{Dokładne podniesienie źródła do oddziaływania kwantowego}

Dla $i<j$ niech $X_{ij}=(|i\rangle\langle j|+|j\rangle\langle i|)/\sqrt2$.
Jest to ortonormalna baza rzeczywistych symetrycznych macierzy
o zerowej przekątnej. Oznaczmy przez $S$ zamianę dwóch czynników,
$|\Omega\rangle=\sum_i|ii\rangle$ oraz $D=\sum_i|ii\rangle\langle ii|$.
Wówczas
\begin{equation}\label{eq:V}
 V=\sum_{i<j}X_{ij}\otimes X_{ij}
   =\frac{|\Omega\rangle\langle\Omega|+S-2D}{2},\qquad
 \Tr_2[V(I\otimes\rho)]=T(\rho).
\end{equation}
To dokładna tożsamość także dla zespolonej hermitowskiej $\rho$.

Trzy ortogonalne projektory spektralne to
\begin{align}
 P_\Omega&=|\Omega\rangle\langle\Omega|/n,&v_\Omega&=(n-1)/2,\\
 P_+&=(I+S)/2-D,&v_+&=1/2,\\
 P_-&=I-P_\Omega-P_+,&v_-&=-1/2.
\end{align}
Dla $n\ge3$ pierwsza przestrzeń własna jest izolowana i jednowymiarowa.
Pasmo $P_-$ obejmuje zarówno przestrzeń antysymetryczną, jak i
bezśladowe wektory diagonalne; nie są to te same przestrzenie.
Ponadto $\Tr_1V=\Tr_2V=0$ oraz
\begin{equation}\label{eq:V2}
 V^2=\frac{I+(n-2)|\Omega\rangle\langle\Omega|}{4}.
\end{equation}

\section{Główny wynik: pełny rząd wymusza korelacje}

\begin{theorem}[Brak nietrywialnych lokalnych pól w komutancie]\label{thm:local}
Dla $n\ge3$ i hermitowskich $h,k$,
\[
 [V,h\otimes I+I\otimes k]=0
 \quad\Longleftrightarrow\quad h,k\text{ są skalarne}.
\]
\end{theorem}

\begin{proof}
Izolowana przestrzeń $\Omega$ wymusza $h+k^T=cI$.
Po utożsamieniu wektora dwuciałowego z macierzą $M$, działanie
$h\otimes I+I\otimes k$ ma postać $cM+[h,M]$.
Dla bezśladowej diagonalnej $M$ składowa pochodząca od urojonej
antysymetrycznej części $h$ jest symetryczna, pozadiagonalna
i należy do pasma $+1/2$, podczas gdy $M$ należy do pasma $-1/2$.
Musi zniknąć dla każdej takiej $M$. Zatem $h$ jest rzeczywista
symetryczna, a $k=cI-h$.

Dla każdej rzeczywistej symetrycznej pozadiagonalnej $M$ komutator
$[h,M]$ jest antysymetryczny, więc przechodzi z pasma $+1/2$ do $-1/2$.
Także on musi zniknąć. Komutacja z każdym $E_{ij}+E_{ji}$ wymusza
$h$ skalarne: trzeci indeks usuwa wszystkie elementy pozadiagonalne,
a porównanie elementów $ij$ zrównuje przekątną. Skalarność $k$ wynika
z wcześniejszej relacji. Kierunek odwrotny jest natychmiastowy.
\end{proof}

\begin{theorem}[Dowolny skończony graf i dowolne pola lokalne]\label{thm:product}
Niech $G$ będzie skończonym połączonym grafem co najmniej dwóch
nośników o wymiarze $n\ge3$. Dla
\[
 H=\sum_a L_a^{(a)}+\sum_{\{a,b\}\in E(G)}g_{ab}V_{ab},\quad g_{ab}\ne0,
\]
gdzie $L_a$ są dowolne hermitowskie, pełnorządowy produkt
$R=\bigotimes_a C_a$ jest stacjonarny wtedy i tylko wtedy,
gdy każde $C_a=I/n$.
\end{theorem}

\begin{proof}
Z $[H,R]=0$ i dodatniości $R$ wynika przez rachunek funkcyjny
$[H,\log R]=0$, gdzie $\log R=\sum_a(\log C_a)^{(a)}$.
Każdy komutator krawędziowy ma zerowe obie jednociałowe ślady
częściowe. Wynika to z $\Tr_1V=\Tr_2V=0$ i cykliczności śladu
częściowego względem operatora działającego na śledzonym czynniku.

Rzut całkowitego komutatora na dwuciałową składową o zerowych
śladach częściowych dla wybranej krawędzi izoluje jej komutator.
Pola lokalne i inne krawędzie nie mogą go skasować. Twierdzenie
\ref{thm:local} wymusza skalarność obu logarytmów na każdej krawędzi.
Spójność grafu i normalizacja stanów kończą dowód. Produkt
maksymalnie mieszany oczywiście komutuje z każdym $H$.
\end{proof}

To jest właściwa treść konieczności korelacji w badanej klasie.
Nie znaczy ona, że każda równowaga musi być splątana.
Pełny rząd jest istotny: produkt $|+_{ij}\rangle\otimes|-_{ij}\rangle$
z równymi rzeczywistymi amplitudami na dwóch etykietach należy
do pasma $-1/2$. Grafy dwudzielne z zerowymi lub zgodnymi polami
lokalnymi dopuszczają naprzemienne produkty tego typu.
Również $n=2$ jest wyjątkiem: $V=X\otimes X/2$ ma nietrywialny
lokalny komutant. Żaden z tych wyjątków nie jest pełnorządowym strict.

\section{Ścisła różnica między polem średnim a skończonym układem}

Dla jawnie przyjętego skalowania
\[
 H_N=-\frac g{N-1}\sum_{a<b}V_{ab},\qquad g=1/\gamma,
\]
oraz początkowego $C^{\otimes N}$ otrzymujemy równanie pola średniego
$\dot C_t=ig[T(C_t),C_t]$. Nie utożsamiamy niezależnych kopii
z fraktalnym uszczegółowieniem nośnika. Zależność granicy od
początkowej hierarchii stanów jest istotna także w ogólnej literaturze
o granicy pola średniego.
\footnote{Z.~Ammari, M.~Falconi, B.~Pawilowski,
\emph{On the rate of convergence for the mean field approximation of
many-body quantum dynamics}, \href{https://arxiv.org/abs/1411.6284}{arXiv:1411.6284}.
Ich ogólnych twierdzeń bosonowych nie przenosi się tu automatycznie
na dowolny produkt stanów mieszanych; poniższy krótko-czasowy szacunek
wynika bezpośrednio ze skończonej hierarchii.}

Dla $f_k(t)=\|\rho_N^{(k)}(t)-C_t^{\otimes k}\|_1$,
$v=\|V\|=(n-1)/2$, hierarchia BBGKY i wzór Duhamela dają
\[
 f_k(t)\le\frac{3gv,k(k-1)}{N-1}t
       +2gvk\int_0^t f_{k+1}(r)\,dr,
\]
przy braku ostatniego składnika dla $k=N$. Jeden składnik
niejednorodny pochodzi od oddziaływań wewnętrznych, dwa od różnicy
współczynników $(N-k)/(N-1)$ i $1$; użyto
$\|[V,R]\|_1\le2v\|R\|_1$. Iteracja do poziomu $N$ daje
\begin{equation}\label{eq:mf}
 f_1(t)\le\frac{3}{2(N-1)}\sum_{r=1}^{N-1}r q^{r+1}
 \le\frac{3q^2}{2(N-1)(1-q)^2},\qquad q=2gvt<1.
\end{equation}
Jest to kontrola na zadanym krótkim przedziale, nie twierdzenie
o dowolnie długim czasie.

Jeżeli $[T(C),C]=0$, to dokładnie
\begin{equation}\label{eq:acceleration}
 \left.\frac{d^2\rho_N^{(1)}}{dt^2}\right|_0
 =-\frac{g^2}{N-1}\Tr_2[V,[V,C\otimes C]].
\end{equation}
Składniki z trzecim nośnikiem znikają dzięki równowadze pola średniego.
Stąd nawet dokładna równowaga Hartreego nie oznacza dokładnej
równowagi skończonego produktu.

Dla rzeczywistego symetrycznego $C$ o stałych przekątnych $C$ i $C^2$
otrzymujemy ponadto zamknięty wzór dwuciałowy, $\tau=gt$:
\begin{align}\label{eq:closed}
 C(\tau)&=C+A(\tau)(C-I/n)
       +B(\tau)(C^2-\Tr(C^2)I/n),\\
 A(\tau)&=\frac{2(\cos\tau-1)}n,\qquad
 B(\tau)=\frac2n\left[\cos\bigl((n/2-1)\tau\bigr)-\cos\tau\right].
\end{align}
Dowód polega na oddzieleniu komutującej zamiany $S$ i rozwinięciu
pozostałej unitarnej części
$I+(e^{-i\tau}-1)D+(e^{i(n/2-1)\tau}-e^{-i\tau})P_\Omega$.
Ślad częściowy daje \eqref{eq:closed} bez dopasowania numerycznego.

Strict $C=I/12+\gamma W$ spełnia założenia. Dla $n=12$ mamy
$C(\pi)=2C/3+I/36$ oraz $C(2\pi)=C$.
Populacje węzłów są stale równe $1/12$, ale widmo lokalnego stanu
zmienia się. Zachowanie entropii globalnej i nierówność subaddytywności
dają $S(C(t))\ge S(C)$ dla produktu początkowego. Nie wynika z tego
monotoniczność: dokładny powrót w $2\pi$ ją wyklucza.

\section{Skorelowana stacjonarność istnieje, ale nie wystarcza}

Niech $A_-=(I-S)/2$, $n\ge3$ i
$\lambda_{\min}(C)\ge1/[2(n-1)]$. Wówczas
\begin{equation}\label{eq:anti}
 R_A(C)=\frac2{n-2}A_-
 \left(C\otimes I+I\otimes C-\frac I{n-1}\right)A_-
\end{equation}
jest stanem, ma oba marginały równe $C$ i komutuje z $V$.
W bazie własnej $C$ jest mieszaniną znormalizowanych projektorów
antysymetrycznych z wagami
$2(c_i+c_j-1/(n-1))/(n-2)$. Wagi są nieujemne, sumują się do jedności,
a suma wag incydentnych dla $i$ wynosi $2c_i$; dowodzi to także
równości marginałów. Stacjonarność wynika z płaskości pasma.

Ta konkretna konstrukcja jest splątana, ponieważ
\[
 \langle\Omega/\sqrt n|R_A(C)^{T_2}|\Omega/\sqrt n\rangle=-1/n.
\]
Transpozycja częściowa każdego stanu separowalnego jest dodatnia,
więc ujemna wartość jest wystarczającym świadkiem. Negatywność jest
co najmniej $1/n$. Nie przypisuje to nośnikom statystyki fermionowej
ani nie dowodzi konieczności splątania dla wszystkich innych uzupełnień.

Dla strict przy $\gamma=1/20$ dokładny przedział spektralny daje
$\lambda_{\min}(C)>1/22$, zatem konstrukcja jest dopuszczalna.
Jest jednak zależna od zadanej odpowiedzi $C$. Co ważniejsze,
każdy wspólny impuls $U\otimes U$ pozostawia ją w tym samym płaskim
paśmie: $R_A(UCU^*)$ także pozostaje stacjonarny, nawet gdy równanie
Hartreego poruszałoby $UCU^*$.
Zapisanie macierzy jądra w stanie nie jest odtworzeniem jej propagatora.
Dodanie osobnego hamiltonianu $W$ byłoby dodatkowym źródłem dynamiki.

\section{Kwantowe odtworzenie kanału z dostarczonego programu}

Istnieje inny, jawnie operacyjny model. Przygotowujemy nieznany stan
sondy $\sigma$ i $N$ niezależnych kopii programu $C$, niezależnych
od sondy. W każdym kroku przykładamy $e^{ig\delta V}$, odrzucamy
użyty program i bierzemy $\delta=t/N$. Kanał jednego kroku to
\[
 \Phi_\delta(\sigma)=\Tr_2\!left[e^{ig\delta V}
                    (\sigma\otimes C)e^{-ig\delta V}\right].
\]
Jest to wariant znanej symulacji hamiltonianu na podstawie próbek
stanu, a nie nowa ogólna zasada kwantowa.
\footnote{S.~Lloyd, M.~Mohseni, P.~Rebentrost,
\emph{Quantum principal component analysis}, Nature Physics 10,
631--633 (2014), \href{https://doi.org/10.1038/nphys3029}{doi:10.1038/nphys3029}.
S.~Kimmel, C.~Y.-Y.~Lin, G.~H.~Low, M.~Ozols, T.~J.~Yoder,
\emph{Hamiltonian Simulation with Optimal Sample Complexity},
\href{https://arxiv.org/abs/1608.00281}{arXiv:1608.00281}, §2 i dodatek B.
Ich ogólnych twierdzeń o optymalności nie utożsamiamy z optymalnością
poniższej szczególnej architektury.}

Dla $h=T(C)$ i $V_0=V-h\otimes I$ zachodzi
\begin{align}
 B_C&=\Tr_2[V_0^2(I\otimes C)]
     =\frac{I+(n-2)C^T}{4}-h^2\succeq0,\\
 \mathcal E_C(\sigma)&=\Tr_2[V_0(\sigma\otimes C)V_0],\qquad
 \mathcal D_C(\sigma)=\mathcal E_C(\sigma)-\tfrac12\{B_C,\sigma\}.
\end{align}
Mapa $\mathcal E_C$ jest całkowicie dodatnia i
$\mathcal E_C^*(I)=B_C$, stąd
$\|\mathcal D_C\|_\diamond\le2\|B_C\|$.
Drugi rząd rozwinięcia jest dokładnie
$\Phi_\delta-\mathcal U_h(\delta)=g^2\delta^2\mathcal D_C+O(\delta^3)$,
gdzie $\mathcal U_h(t)$ jest sprzężeniem przez $e^{igth}$.

Trzeci rząd reszty kanału unitarnego ma normę co najwyżej
$\delta^3(2\|H\|)^3/6$. Teleskopowanie kanałów, których normy
diamentowe są równe jeden, daje ścisłą granicę
\begin{equation}\label{eq:diamond}
 \|\Phi_{t/N}^N-\mathcal U_h(t)\|_\diamond
 \le\frac{2g^2t^2\|B_C\|}{N}
 +\frac{4g^3|t|^3}{3N^2}\bigl(v^3+\|h\|^3\bigr),
\end{equation}
ograniczoną dodatkowo przez $2$. Dotyczy ona także sondy splątanej
z rejestrem odniesienia.

Przy $C=I/12+W/20$, $g=20$, dokładnie opłacone $\|W\|<3$ daje
\begin{equation}\label{eq:resource}
 \|\Phi_{t/N}^N-\mathcal U_W(t)\|_\diamond
 \le\min\left(2,\frac{2000t^2}{3N}
                   +\frac{5324108|t|^3}{3N^2}\right).
\end{equation}
Dla $t=1$, $N=10^6$ prawa strona wynosi
$501331027/750000000000<7\cdot10^{-4}$.
To oszacowanie teoretyczne, nie wynik aparatury.
Kanał $\mathcal U_W$ jest równy kanałowi laplasjanowemu strict
z dokładnością do globalnej fazy propagatora. Nie wyprowadzono tu
kontrolowanej bramki $U$, zera energii ani absolutnej fazy zegara.

\section{Nieusuwalny błąd w zadanej architekturze}

Dla jednorodnego wektora Perrona $u$ połóżmy $P=|u\rangle\langle u|$.
Każdy program $\rho$ spełniający $T(\rho)=\gamma W$ ma tę samą liczbę
$c_0=\langle u|\rho|u\rangle=1/n+\gamma s$.
Bez względu na ukryte dane przekątnej i koherencje urojone programu,
dokładna utrata prawdopodobieństwa pozostania sondy w $P$ po jednym
zderzeniu o fazie $\tau=g\delta$ wynosi
\begin{equation}\label{eq:loss}
 L(\tau,c_0)=(1-c_0)\left(1-\frac4{n^2}\right)\sin^2\frac\tau2
     +\frac{4(n-1)c_0}{n^2}\sin^2\frac{(n-2)\tau}{4}.
\end{equation}
Aby to sprawdzić, wystarczy obliczyć operator
$\langle u|e^{i\tau V}|u\rangle$ na programie. Ma on postać
$a_\tau I+b_\tau P$; kwadraty jego wartości na $u^\perp$ i $u$
dają oba składniki \eqref{eq:loss}.

Współczynnik $\tau^2$ to
\[
 f_0=\frac{n^2-4}{4n^2}+\frac{(n-2)(n-4)}{4n}c_0.
\]
Dla $n=12$ jest dodatni. Ponieważ $P$ komutuje z docelowym $W$,
rozwinięcie $\log\Phi_\delta$ przy ustalonym programie i $t$ daje
\begin{equation}\label{eq:lower}
 1-\Tr[P\Phi_{t/N}^N(P)]
     =\frac{g^2t^2f_0}{N}+O(N^{-2}).
\end{equation}
Jest to dolny świadek odległości śladowej, nie tylko norma pomocnicza.
Dla $\gamma=1/20$ współczynnik $g^2f_0$ wynosi około $208.12135$.
Dokładne potęgowanie kanału daje przy $N=10^6$ około $208.09761$
dla $N$ razy utrata; nie traktujemy tej liczby zmiennoprzecinkowej
jako dowodu asymptotyki.

Równanie \eqref{eq:lower} nie jest dolną granicą dla wszystkich
algorytmów kwantowych, programów skorelowanych ani dowolnych reguł
sterowania. Kopie i architektura są częścią twierdzenia.
Pełnorządowe programy podlegają również standardowemu ograniczeniu
programowalności.
\footnote{M.~A.~Nielsen, I.~L.~Chuang, \emph{Programmable Quantum Gate Arrays},
Physical Review Letters 79, 321 (1997),
\href{https://arxiv.org/abs/quant-ph/9703032}{arXiv:quant-ph/9703032}.}
Jeżeli jedno ustalone urządzenie daje dokładny kanał unitarny dla
pełnorządowego programu, to przez rozkład tego programu jako
$\epsilon\rho+(1-\epsilon)\tau$ i ekstremalność kanałów unitarnych
musi dawać ten sam kanał dla każdego $\rho$.
Skończony pełnorządowy program nie może zatem dokładnie wybierać
różnych unitariów w takim urządzeniu.

\section{Optymalizacja programu nie jest źródłem jądra}

Twirl po grupie dihedralnej przekształca dowolny dopuszczalny program
w $C_\gamma=I/n+\gamma W$, zachowując źródło $T$.
Jest to mieszanina zwykłych unitarnych permutacji, nie operacja
antyunitarna. Dlatego jakikolwiek program istnieje wtedy i tylko wtedy,
gdy
\[
 0<\gamma\le\gamma_*:=1/(n\ell).
\]
Kowariancja $B_C$ i wypukłość jego największej wartości własnej
pokazują, że $C_\gamma$ minimalizuje $\|B_C\|$ przy ustalonym
$\gamma$. Nie dowodzimy jednoznaczności wszystkich minimizerów.

Dla $n=12$ wartości własne $B_C/\gamma^2$ mają postać
$11/(24\gamma^2)+5\lambda/(2\gamma)-\lambda^2$.
Są rosnące w $\lambda$ na całym dopuszczalnym zakresie, gdyż
$5/(2\gamma)-2s\ge30\ell-2s\ge8\ell>0$.
Wartość Perrona maleje z $\gamma$, więc optimum tego oszacowania
osiąga się przy $\gamma_*$:
\begin{align}
 \min\frac{\|B_C\|}{\gamma^2}&=66\ell^2+30\ell s-s^2,\\
 \min\frac{f_0}{\gamma^2}&=55\ell^2+20\ell s.
\end{align}
Dokładne zewnętrzne przedziały Fourierowskie dają odpowiednio około
$\gamma_*=0.1222120805$, $61.8939393$ i $48.2148582$.
Są to warunkowe parametry kodowania i błędu, nie wyprowadzone stałe fizyczne.

Sama wariancja $B_C$ nie jest uniwersalnym kosztem operacyjnym:
oddziaływanie $I\otimes Z$ może mieć dodatnią wariancję i jednocześnie
nie zmieniać sondy. Dlatego ważny jest niezależny świadek
\eqref{eq:loss}, a nie tylko optymalizacja górnej granicy.
Ponadto dwa programy o tym samym $T(\rho)$ mogą mieć różny szum;
dla jego generatora drugiego rzędu
$\mathcal D_\rho(I)=(n-2)(\rho-\rho^T)/4$.
Samo $W$ nie ustala całej dynamiki poprawki kwantowej.

\section{Granica arytmetyczna i osobna klasyfikacja legacy}

Wcześniejszy P512/O216 dowodzi, że stosunek dwóch niezerowych
częstotliwości laplasjanowych strict, $\lambda_2/\lambda_1$, jest
transcendentalny dla dokładnych parametrów \eqref{eq:strict}.
Nie jest to nowe twierdzenie tego raportu.
\footnote{Źródło repozytoryjne:
\texttt{FIN\_Programs\_507\_516\_Nonlinear\_Source\_and\_Certified\_Localization\_Report\_EN.tex},
§P512; \texttt{fin\_programs\_507\_516\_research.py}, funkcja
\texttt{p512\_frequency\_arithmetic}. Dowód używa
$z=e^{i/4000}$ i nieproporcjonalnych wielomianów Laurenta.}

Nowe zastosowanie jest następujące. Każdy skończony $H_N$ powyżej
ma wymierne elementy z dokładnością do jednej skali, więc jego
stosunki niezerowych częstości Bohra są algebraiczne.
Żadne ustalone liniowe przygotowanie i odczyt nie odtworzą zatem
dokładnie całej krzywej powrotu $|[e^{-itA}]_{00}|^2$ strict.
Krzywa ta zawiera obie częstotliwości $\lambda_1,\lambda_2$ z
dodatnimi współczynnikami; jednoznaczność skończonych sum wykładniczych
wymuszałaby ich obecność w widmie mikroskopowym, co daje sprzeczność.
Nie wyklucza to przybliżenia, nieskończonej granicy, dodatkowych
niezależnie wyprowadzonych sprzężeń ani innego sterowania czasowego.
Nie jest to również dolna granica błędu odporna na dowolną zmianę
dokładnych parametrów wewnątrz niepewności pomiarowej.

Dla osobnego kanonicznego cyklu legacy
$W_d=4\ln2\cos(\pi d/4+\pi/6)/(1+0.01d)$ stosunki widmowe są
algebraiczne: wspólny czynnik $4\ln2$ się skraca, a reszta wag
i współczynniki Fouriera są algebraiczne. Przeszkody P512 nie
wolno na niego przenieść. Nie dowodzi to realizacji legacy przez $V$.
Twierdzenie o produktach oraz konstrukcja \eqref{eq:anti} stosują się
natomiast osobno do obu referencji. Dla legacy i $\gamma=1/1000$
oszacowanie $|W_{ij}|<3$ daje $C\succeq(1/12-33/1000)I\succ I/22$.
Ujemne wagi legacy nie są tu interpretowane jako dodatnie tempo przejść.

\section{Bilans naukowy i następny test rozstrzygający}

\paragraph{Udowodniono.}
Podniesienie \eqref{eq:V} ma trywialny lokalny komutant, co wyklucza
niejednorodne pełnorządowe produkty stacjonarne na dowolnym skończonym
połączonym grafie. Strict wymaga w tej klasie korelacji; jego równowaga
pola średniego nie jest dokładną równowagą produktu mikroskopowego.
Jednocześnie istnieją jawne skorelowane uzupełnienia oraz kontrolowany
protokół kwantowego odtworzenia kanału, ze skończonym budżetem kopii,
górną granicą błędu i dolnym świadkiem w zadanej architekturze.

\paragraph{Obalono w zadanych klasach.}
Nie wystarcza potraktowanie stacjonarności pola średniego jako
stacjonarności niezależnych składników. Nie wystarcza znalezienie
skorelowanego stanu o prawidłowym marginałe: pasmo płaskie pokazuje
różnicę między pamięcią a propagacją. Nie można też usunąć skończonego
szumu badanego procesora przez dowolną zmianę danych programu
niewidocznych dla $T$.

\paragraph{Pozostaje otwarte.}
Konkretne prawo mikroskopowe, identyfikacja jego stopni swobody,
przygotowanie korelacji lub programu, dopływ kopii i zegar nie zostały
wyprowadzone z nadsolitonu. Są to rzeczywiste, dopuszczalne zasoby
modelu operacyjnego, lecz nie dowód samoistnego powstania $W$.
Stan początkowy może być legalnym wejściem teorii fizycznej; samo
jego występowanie nie jest argumentem przeciw wszelkiej fizyce FIN.

Najbardziej rozstrzygający dalszy kierunek to wyprowadzenie jednego
konkretnego prawa mikroskopowego i jego stanu korelacyjnego lub
strumienia programu, z niezależnie określoną procedurą pomiaru.
Powinno ono przewidzieć jednocześnie propagację i skończone poprawki,
nie tylko odtworzyć macierz wpisaną do danych wejściowych.
Alternatywa kwantowa wymaga jawnych korelacji i kontroli granicy;
alternatywa klasyczna lub efektywna wymaga własnego prawa operacyjnego.

Nie uzyskano dowodu pierwszeństwa wobec ogólnej literatury, nowej
cząstki, fizycznego źródła jednostek, mostu legacy--strict, transferu
ról, rozstrzygnięcia QW-2191, SM/GR, $L_{\rm total}$ ani ToE.
Nadsoliton jest pierwotną informacją w stanie solitonicznym;
konstrukcje tensorowe nie ustanawiają pod nim dodatkowej warstwy.

\section*{Źródła i dane sprawdzalne}

\begin{enumerate}
\item Dokumenty repozytoryjne: \texttt{The FIN Kernel as an Unknown Mat.md},
§4.1; \texttt{fin\_projected\_learning/REPORT.md}, §28;
\texttt{fin\_transition\_hierarchy/REPORT.md}; wcześniejszy raport
\texttt{FIN\_Adaptive\_Feedback\_Operational\_Obstruction.tex}.
\item Z.~Ammari, M.~Falconi, B.~Pawilowski,
\emph{On the rate of convergence for the mean field approximation of
many-body quantum dynamics}, \href{https://arxiv.org/abs/1411.6284}{arXiv:1411.6284}.
\item S.~Lloyd, M.~Mohseni, P.~Rebentrost,
\emph{Quantum principal component analysis}, Nature Physics 10,
631--633 (2014), \href{https://doi.org/10.1038/nphys3029}{doi:10.1038/nphys3029}.
\item S.~Kimmel, C.~Y.-Y.~Lin, G.~H.~Low, M.~Ozols, T.~J.~Yoder,
\emph{Hamiltonian Simulation with Optimal Sample Complexity},
\href{https://arxiv.org/abs/1608.00281}{arXiv:1608.00281}.
\item M.~A.~Nielsen, I.~L.~Chuang,
\emph{Programmable Quantum Gate Arrays}, Physical Review Letters 79,
321 (1997), \href{https://arxiv.org/abs/quant-ph/9703032}{arXiv:quant-ph/9703032}.
\item P512/O216: wskazany w przypisie raport repozytoryjny i jego
certyfikat stosunku częstotliwości. Dokładne przedziały wag i widma:
\texttt{fin\_replication\_consistency/certify.py} oraz
\texttt{fin\_projected\_learning/research.py}.
\item Dane tego opracowania: \texttt{fin\_quantum\_lift/results.json},
\texttt{collision\_results.json} i odpowiadające im implementacje
\texttt{research.py}, \texttt{collisions.py} oraz testy.
Przedziały wymierne, tożsamości algebraiczne i wyniki zmiennoprzecinkowe
mają osobne statusy; żadna mała reszta nie zastępuje dowodu ogólnego.
\end{enumerate}
\end{document}
